\section{Federated control as a research direction}
\label{sec:federation}

The current demonstrator uses one operator process to dispatch work to Docker engines; the
network optimiser likewise observes all sites at once. This is a useful benchmark, but it assumes
that its controller can always reach the sites. A residential fleet may instead experience
unreliable home connections, changing heat demand, and contractual restrictions on where data may
run. We therefore propose a \emph{federated} control plane as future work, not as a feature or
measurement of the present demonstrator (decision 0035).
The repository contains a deterministic claim-validation model, but no distributed ownership
service or real-node federation.

\paragraph{Two time scales.} Heat control stays local: each site measures demand and tank state and
applies the same heat-led scheduler and executor interface used in the annual model. It can reduce
or pause interruptible work when disconnected. Job placement is a separate, slower agreement:
sites publish expiring offers of free compute resources and useful-heat headroom. The job origin
filters offers by the customer's locality policy and job requirements, then compares eligible
sites. The heat signal helps rank offers but cannot override a data-sovereignty restriction.

\paragraph{The ownership problem.} An offer cannot itself authorize execution. Under a network
partition two sites may otherwise accept the same job. We propose a durable claim containing the
job identity, attempt generation, owner, policy revision, and expiration, agreed by a replicated
quorum and checked by the node executor. Old generations are fenced. Reassignment waits for
confirmed old-process termination or lease expiry plus a bounded clock margin: fencing alone
cannot stop an isolated process. An uncertain Docker start is reconciled before a retry. This
deliberately sacrifices new-job availability in a minority
partition to prevent duplicate execution. Running jobs may continue only within their valid
authorization and local heat limits. Non-interruptible work needs a declared run window; it does
not gain automatic pause semantics. A logically unique ownership decision remains necessary even
without one permanently designated coordinator.

\paragraph{Hypotheses to test.} First, node-local heat control should preserve useful heat when
the placement service is unavailable, while a controller-dependent implementation cannot issue
new commands. Second, federated offers may reduce missed deadlines when the coordinator is
unreachable, provided a quorum and eligible sites remain reachable. Third, stale heat offers can
increase dumping or backup heat; offer age should be measured as an explicit input, not hidden in
the routing score. These are hypotheses, not findings.

\paragraph{Experiment design.} Use the same weather, job arrivals, locality policies, heat
balance and economic parameters for the existing central B2/LP baselines and a federated variant.
Inject reproducible delays, partitions, node crashes, stale offers and clock uncertainty; report
useful heat, dumped heat, missed deadlines, sold work, duplicate attempts, unavailable placement
time and control traffic. A two-device demonstration is insufficient to test quorum behavior:
three independent voting failure domains are needed to survive one partitioned voter while
retaining majority. Real-device results must separately document the hardware, clock assumptions,
failure injection, command and logs. No numerical advantage is claimed until that experiment is
run.
